\chapter*{Introduction générale}
\addcontentsline{toc}{chapter}{Introduction générale}
\markboth{Introduction générale}{}

Pour un établissement bancaire, la sécurité du code source est aujourd'hui un enjeu
de premier plan. La généralisation des pratiques DevOps a accéléré le rythme des
livraisons, mais elle a aussi élargi la surface d'attaque que chaque déploiement
expose. Une vulnérabilité qui franchit l'étape de fusion sans avoir été repérée peut
compromettre l'intégrité du système d'information et engager la responsabilité de la
banque, tant vis-à-vis de ses clients que de ses régulateurs. Les exemples ne manquent
pas : une injection SQL au cœur d'une API de paiement, un secret d'authentification
commité par mégarde dans un dépôt Git, ou une image Docker bâtie sur une librairie
exposée à un CVE critique. Or la revue de sécurité reste, dans le modèle classique,
un travail d'experts qui relisent chaque \textit{Pull Request} à la main. Le procédé
est lent, son résultat dépend largement du relecteur, et il passe d'autant plus mal
à l'échelle que les équipes de développement s'agrandissent.

Deux évolutions récentes rendent une autre approche envisageable. Les grands modèles
de langage (LLM) savent désormais analyser du code source et en formuler des
corrections lisibles. De leur côté, les outils d'analyse de sécurité statique (SAST)
ont gagné en maturité au point d'être exploitables en production. En les combinant,
il devient possible de déclencher une revue automatisée à chaque \textit{Pull Request} :
en quelques minutes, celle-ci produit une analyse alignée sur l'OWASP Top 10 et publie
ses commentaires de correction directement sur GitHub. C'est précisément l'objet de ce
projet de fin d'études, réalisé au sein de la Banque de Tunisie et des Emirats (BTE).
Il s'agit de concevoir et de déployer le \textbf{BTE Security AI Agent}, un moteur
DevSecOps autonome piloté par une intelligence artificielle exécutée localement.
L'agent intercepte les événements de \textit{Pull Request} GitHub, orchestre en
parallèle cinq outils d'analyse de sécurité et produit une revue complète en 15 à
25 minutes selon la complexité de la \textit{Pull Request}, sans qu'aucun
code source ne quitte l'infrastructure interne de la banque.

Le rapport s'organise en quatre chapitres. Le \textbf{premier chapitre} présente
l'organisme d'accueil et le cadre général du projet, pose la problématique et expose
la solution retenue ; il revient également sur la méthodologie Scrumban et sur le
déroulement du travail, réparti en dix sprints sur cinq mois. Le \textbf{deuxième
chapitre} pose les bases théoriques (DevSecOps, inférence locale par grands modèles
de langage, outils d'analyse statique), puis justifie les technologies retenues et
spécifie les besoins fonctionnels et non fonctionnels. Le \textbf{troisième chapitre}
détaille l'architecture logicielle : le workflow LangGraph à neuf nœuds, l'architecture
à deux modèles LLM, le schéma de la base de données et l'assistant de chat opérationnel
doté de son système anti-hallucination. Le \textbf{quatrième chapitre} décrit
l'environnement de déploiement et le travail mené phase par phase, de la configuration
des services Docker à la mise en place de l'observabilité, en passant par l'intégration
des outils de sécurité et l'implémentation du pipeline LLM. Une conclusion générale
referme le document : elle récapitule le travail accompli et trace les perspectives
d'amélioration.
